\documentclass[11pt]{article}
\usepackage[margin=1in]{geometry}
\usepackage{amsmath}
\usepackage{booktabs}
\usepackage{hyperref}
\usepackage{longtable}

\title{War Carbon Emissions Dashboard \\ Methodology v1.0}
\author{WCED Scientific Steering Committee}
\date{2026-05-23}

\begin{document}
\maketitle

\section{Scope}

This document specifies methodology version \texttt{v1.0} of the War Carbon
Emissions Dashboard (WCED). V1 quantifies $\mathrm{CO_2}$ released by the
combustion of hydrocarbons at oil and fuel infrastructure (refineries, crude
and product depots, petrochemical complexes) inside the Iran--Gulf theatre of
the 2026 Iran--Israel--US war, from 2026-02-28 onward. It does \emph{not}
cover destroyed-building embodied carbon, munitions embodied carbon, combat
aviation fuel, shipping rerouting, reconstruction, or non-carbon ecological
damage. Each of those is deferred to a later version (see
\texttt{docs/V1\_PLAN.md}, Section ``How V1 Evolves''). Numbers in V1 are
estimates of $\mathrm{CO_2}$ released from observed fires; they are not
attributed to specific belligerents.

\section{Data Model}

Every quantitative output crosses module boundaries as one of three Pydantic
v2 types:

\begin{description}
\item[\texttt{Distribution}] Carries \texttt{p5}, \texttt{p50},
  \texttt{p95}, the underlying Monte Carlo \texttt{samples}, and a
  \texttt{provenance\_id}. Emission estimates are \emph{never} returned as
  scalars. Scalar arithmetic on a \texttt{Distribution} uses
  \texttt{apply\_scalar(factor, provenance\_id=factor\_record\_id)} so that
  the scalar's own source is recorded; silent inheritance of the parent
  provenance is forbidden.
\item[\texttt{ProvenanceRecord}] Records \texttt{produced\_by},
  \texttt{inputs} (list of upstream provenance IDs), \texttt{method} (versioned
  string), \texttt{parameters} (free-form), \texttt{produced\_at}, and an
  optional \texttt{confidence\_label}. Every number in the database is
  reachable from a chain of these records terminating in a primary source.
\item[\texttt{ConfidenceLabel}] One of
  \texttt{CONFIRMED}, \texttt{VERIFIED}, \texttt{REPORTED}, \texttt{SUSPECTED},
  \texttt{CLAIMED}. Assigned by the verification pipeline (\S4) and
  propagated forward into every downstream estimate.
\end{description}

\section{Emission Calculations}

V1 produces two independent emission estimates per fire event --- a top-down
FRP-based estimate (\S3.3) and a bottom-up inventory-based estimate
(\S3.4) --- and reconciles them (\S3.5). All scalar coefficients are loaded
from \texttt{data/emission\_factors.yaml}; all event/facility-level priors
are loaded from \texttt{data/parameter\_distributions.yaml}.

\subsection*{\S3.1--3.2 (reserved)}
Sections 3.1 (facility registry) and 3.2 (factor tables) are documented in
the YAML files themselves and not duplicated here.

\subsection{FRP-based method}
\label{sec:frp}

Fire Radiative Power (FRP) observations from VIIRS, MODIS, and Sentinel-3
SLSTR are clustered into an event-level time series. The \emph{raw} Fire
Radiative Energy is the trapezoidal integral of the observed overpass
samples:
\begin{equation}
I_{\mathrm{raw}} \;=\; \int_{t_0}^{t_1} \mathrm{FRP}(t)\,\mathrm{d}t.
\label{eq:fre-raw}
\end{equation}
The \emph{corrected} Fire Radiative Energy $I_{\mathrm{FRE}}$ used by all
downstream equations applies two multiplicative corrections:
\begin{equation}
I_{\mathrm{FRE}} \;=\; k_{\mathrm{ext}} \cdot d \cdot I_{\mathrm{raw}},
\label{eq:fre}
\end{equation}
where $d$ is the burn duty cycle (\texttt{burn\_duty\_cycle},
triangular(0.4, 0.7, 0.95)) accounting for fire activity in unobserved
intervals between overpasses, and $k_{\mathrm{ext}}$ is the FRP
extrapolation factor (\texttt{frp\_extrapolation\_factor},
$\mathcal{N}(1.0, 0.15)$) correcting residual bias when the overpass
sampling is sparse. Both $d$ and $k_{\mathrm{ext}}$ are sampled
independently per event in the Monte Carlo (\S5); they must \emph{never}
be pre-baked into $I_{\mathrm{raw}}$, since collapsing them to point
values would understate the p5--p95 width of every FRP-based estimate.
$I_{\mathrm{raw}}$ and $I_{\mathrm{FRE}}$ both have units of MJ.

The fuel mass combusted and resulting $\mathrm{CO_2}$ mass are then
\begin{align}
m_{\mathrm{fuel}} &= \alpha \cdot I_{\mathrm{FRE}},
\label{eq:fuel}\\
m_{\mathrm{CO_2}} &= m_{\mathrm{fuel}} \cdot \beta \cdot r,
\label{eq:co2-frp}
\end{align}
with $\alpha$ the FRP-to-combustion-rate coefficient
(\texttt{frp\_to\_combustion\_rate}, $\mathcal{N}(0.368, 0.05)$ kg/MJ;
Wooster et al.\ 2005), $\beta = 44/12 \cdot f_C$ the stoichiometric
$\mathrm{C}\!\to\!\mathrm{CO_2}$ conversion with $f_C \approx 0.86$ the
carbon mass fraction of crude (AP-42), and $r$ the carbon recovery as
$\mathrm{CO_2}$ (\texttt{carbon\_recovery\_as\_co2},
triangular(0.92, 0.96, 0.98); Hobbs \& Radke 1992). The product
$\beta \cdot r$ collapses to the familiar $\approx 3.15$
$\mathrm{kg\,CO_2}$ per kg crude burned used by the CCI Iran brief
when central values are substituted.

\subsection{Inventory-based method}
\label{sec:inventory}

For facilities with a curated capacity $C$ (barrels) in the
\texttt{data/facilities/} registry, the bottom-up estimate is
\begin{equation}
m_{\mathrm{CO_2}} \;=\; C \cdot \phi \cdot \psi \cdot \mathrm{EF},
\label{eq:co2-inv}
\end{equation}
where $\phi$ is the fraction of nameplate capacity present at strike
(\texttt{facility\_inventory\_at\_strike}, uniform(0.3, 0.9)), $\psi$ is
the destroyed fraction observed from before/after Sentinel-2 imagery (event-
specific; default triangular($\hat\psi - 0.15$, $\hat\psi$, $\hat\psi + 0.15$)
truncated to $[0, 1]$), and $\mathrm{EF}$ is the appropriate combustion
emission factor:

\begin{itemize}
\item Crude depots / refinery crude units:
  \texttt{crude\_oil\_combustion} = $0.425$ tCO$_2$/barrel (EPA AP-42
  \S1.3; CCI Iran brief 2026).
\item Refined-product depots: \texttt{refined\_product\_combustion}
  = $0.430$ tCO$_2$/barrel (IPCC 2006 GL Vol.\ 2 Ch.\ 1 Table 1.3).
\end{itemize}

Each factor entry in \texttt{emission\_factors.yaml} carries an
\texttt{applicable\_facility\_types} list. \texttt{quantify/inventory.py}
must validate that the selected factor is applicable to the event's
facility type and raise if not; applying a crude factor to a gas-processing
facility silently produces wrong numbers.

\subsection{Reconciliation}
\label{sec:reconcile}

Both estimates are Monte Carlo distributions. The reconciliation rule is:

\begin{enumerate}
\item Compute $\rho = m_{\mathrm{CO_2}}^{\mathrm{inv,p50}} \,/\,
  m_{\mathrm{CO_2}}^{\mathrm{frp,p50}}$.
\item If $0.5 \le \rho \le 2.0$: the two estimates agree. Publish both,
  flag the event \texttt{reconciled\_ok}, and display the
  \emph{envelope distribution}
  $\mathrm{Distribution.union}(D_{\mathrm{frp}}, D_{\mathrm{inv}})$
  as the headline number (p5/p50/p95 over pooled samples).
\item If $\rho < 0.5$ or $\rho > 2.0$: flag \texttt{reconciliation\_review}.
  Both estimates are stored; neither is promoted to the public dashboard
  until editorial review (\S4.4) resolves the discrepancy. Publish only
  after a written reviewer note enters the changelog.
\item If a third, ``reported'' value exists (CEOBS / CCI / government
  brief), it is recorded as a \texttt{Claimed}-confidence cross-check and
  never enters the headline arithmetic.
\end{enumerate}

\section{Verification and Confidence Assignment}
\label{sec:verify}

Each candidate fire event is independently assessed against three evidence
streams: FIRMS persistence (number of overpasses with a hotspot at the
location), Sentinel-2 optical classification (was a clear-sky scene found,
and did the classifier label it \texttt{CONFIRMED\_FIRE}), and ACLED
corroboration (one or more armed events within the spatio-temporal window).
The assignment is implemented in \texttt{wced/verify/confidence.py} and
reproduced as Table~\ref{tab:confidence}.

\begin{table}[h]
\centering
\caption{Confidence label decision table (v1.0).}
\label{tab:confidence}
\begin{tabular}{lcccl}
\toprule
\textbf{Label} & \textbf{FIRMS persistent ($\ge 2$ passes)} &
\textbf{S2 confirms fire} & \textbf{ACLED match} & \textbf{Notes} \\
\midrule
\texttt{CONFIRMED} & yes & yes & yes & all three streams agree \\
\texttt{VERIFIED}  & yes & yes & no  & satellite-only, no ACLED yet \\
\texttt{REPORTED}  & yes & no  & any & cloud/no-scene or S2 ambiguous \\
\texttt{SUSPECTED} & no  & --  & any & ACLED-only \emph{never} auto-promotes \\
\texttt{SUSPECTED} & no  & --  & no  & single overpass, weakest non-claimed \\
\texttt{CLAIMED}   & --  & --  & --  & news/state source only, no satellite ev. \\
\bottomrule
\end{tabular}
\end{table}

The persistence threshold is fixed at $\ge 2$ overpasses
(\texttt{\_PERSISTENT\_MIN\_OVERPASSES}). The ACLED-only case is
intentionally non-promoting: an armed event near a facility without a
FIRMS hotspot may be a near-miss, a non-incendiary strike, or a
mis-geocoded record. Such candidates are routed to editorial review with
a logged warning, never auto-confirmed.

Only \texttt{CONFIRMED} and \texttt{VERIFIED} events feed the headline
dashboard total. \texttt{REPORTED} events appear on the map with a distinct
marker. \texttt{SUSPECTED} and \texttt{CLAIMED} events stay in the
pending-review queue.

\section{Uncertainty Quantification}
\label{sec:uq}

Every emission estimate is the output of a Monte Carlo simulation with
$N = 10{,}000$ samples. The parameter PDFs are loaded verbatim from
\texttt{emission\_factors.yaml} and \texttt{parameter\_distributions.yaml};
Table~\ref{tab:priors} reproduces the central values used in v1.0.

\begin{table}[h]
\centering
\caption{Monte Carlo parameter priors (methodology v1.0).}
\label{tab:priors}
\begin{tabular}{llll}
\toprule
\textbf{Parameter} & \textbf{Distribution} & \textbf{Units} &
\textbf{Source} \\
\midrule
\texttt{crude\_oil\_combustion}      & triangular(0.405, 0.425, 0.445) &
  tCO$_2$/barrel & EPA AP-42 \S1.3 \\
\texttt{refined\_product\_combustion} & triangular(0.410, 0.430, 0.455) &
  tCO$_2$/barrel & IPCC 2006 GL Vol.\ 2 \\
\texttt{frp\_to\_combustion\_rate}   & $\mathcal{N}(0.368,\,0.05)$ &
  kg/MJ          & Wooster et al.\ 2005 \\
\texttt{carbon\_recovery\_as\_co2}   & triangular(0.92, 0.96, 0.98) &
  --             & Hobbs \& Radke 1992 \\
\texttt{burn\_duty\_cycle}           & triangular(0.40, 0.70, 0.95) &
  --             & expert prior (Kuwait '91; Abqaiq '19) \\
\texttt{facility\_inventory\_at\_strike} & uniform(0.30, 0.90) &
  --             & expert OSINT prior \\
\texttt{frp\_extrapolation\_factor}  & $\mathcal{N}(1.0,\,0.15)$ &
  --             & internal MODIS/VIIRS calibration \\
\bottomrule
\end{tabular}
\end{table}

Sampling rules:
\begin{itemize}
\item Triangular distributions are sampled on $[\text{low}, \text{high}]$
  with mode at \texttt{value}; \texttt{low}/\texttt{high} default to
  \texttt{uncertainty\_low}/\texttt{uncertainty\_high} when explicit fields
  are absent.
\item Normal distributions use the listed \texttt{sigma} directly; the
  documented \texttt{uncertainty\_low}/\texttt{high} are recorded for
  reviewer cross-check only and are \emph{not} used as truncation bounds.
\item Uniform distributions use strict \texttt{low}/\texttt{high} support.
\item Parameter samples are drawn independently per event (no shared
  random state across events; correlation across facilities is a V2
  question).
\item Aggregation across events: per-event \texttt{samples} arrays are
  summed element-wise to yield an aggregate sample array, then
  p5/p50/p95 are computed on the sum.
\end{itemize}

Cache keys for any persisted Monte Carlo output must include the
\texttt{methodology\_version} string. Recomputing all estimates after a
version bump is a deliberate operation, never automatic.

\section{Worked Example: Shahran Depot}
\label{sec:example}

\textbf{Inputs.}
\begin{itemize}
\item Facility: Shahran fuel depot (Tehran), nameplate capacity
  $C = 5.00 \times 10^{5}$ barrels (crude/refined-product blend; treat as
  crude for V1).
\item Observed destroyed fraction $\hat\psi = 0.40$
  (Sentinel-2 before/after; triangular(0.25, 0.40, 0.55)).
\item Raw FRP integral $I_{\mathrm{raw}} = 8.5 \times 10^{7}$~MJ
  ($\approx 85$~TJ) from trapezoidal interpolation of 4 VIIRS/MODIS
  overpasses spanning an 18-hour window
  (mean FRP $\approx 1.3$~GW, consistent with a major depot fire).
\item Confidence label \texttt{CONFIRMED} (FIRMS $n_{\text{passes}}=4$;
  S2 \texttt{CONFIRMED\_FIRE}; 1 ACLED match within $\pm 24$~h).
\end{itemize}

\textbf{FRP-based estimate (point values).} Apply the corrections of
Eq.~\eqref{eq:fre} at prior central values ($k_{\mathrm{ext}}=1.0$,
$d=0.70$), then Eqs.~\eqref{eq:fuel}--\eqref{eq:co2-frp}:
\begin{align*}
I_{\mathrm{FRE}} &= 1.0 \cdot 0.70 \cdot 8.5\times 10^{7}
   = 5.95 \times 10^{7} \text{ MJ}, \\
m_{\mathrm{fuel}} &= 0.368 \cdot 5.95\times 10^{7}
   = 2.19 \times 10^{7} \text{ kg} \;\approx\; 2.19 \times 10^{4} \text{ t fuel}, \\
m_{\mathrm{CO_2}}^{\mathrm{frp}} &= 2.19 \times 10^{7} \cdot 3.15
   \;\approx\; 6.90 \times 10^{7} \text{ kg}
   \;=\; \mathbf{6.90 \times 10^{4} \text{ tCO}_2}.
\end{align*}

\textbf{Inventory-based estimate (point values).} From Eq.~\eqref{eq:co2-inv}
with $\phi = 0.60$ (uniform midpoint) and $\psi = 0.40$:
\begin{align*}
m_{\mathrm{CO_2}}^{\mathrm{inv}} &= 5.00\times 10^{5} \cdot 0.60 \cdot 0.40
  \cdot 0.425
  \;=\; 5.10 \times 10^{4} \text{ tCO}_2.
\end{align*}

\textbf{Reconciliation.} $\rho = 5.10 \times 10^{4} / 6.90 \times 10^{4}
\approx 0.74$, inside the $[0.5, 2.0]$ band, so the event is flagged
\texttt{reconciled\_ok} and the envelope distribution is published.
\emph{Reviewer note:} events where $\rho$ falls within 10\,\% of either
boundary (i.e.\ $\rho \in [0.50, 0.55] \cup [1.82, 2.00]$) should receive
extra editorial scrutiny --- a small shift in the destroyed-fraction prior
or the burn duty cycle could push them into \texttt{reconciliation\_review}.

\textbf{Monte Carlo result ($N = 10{,}000$, methodology v1.0; reproducible
fixture values used as test targets).}

\begin{table}[h]
\centering
\begin{tabular}{lrrr}
\toprule
Estimate & $p_5$ (tCO$_2$) & $p_{50}$ (tCO$_2$) & $p_{95}$ (tCO$_2$) \\
\midrule
FRP-based              & $3.5 \times 10^{4}$ & $6.9 \times 10^{4}$ & $1.15 \times 10^{5}$ \\
Inventory-based        & $2.4 \times 10^{4}$ & $5.1 \times 10^{4}$ & $9.2 \times 10^{4}$ \\
\textbf{Envelope (headline)} &
  $\mathbf{2.6 \times 10^{4}}$ & $\mathbf{6.0 \times 10^{4}}$ & $\mathbf{1.10 \times 10^{5}}$ \\
\bottomrule
\end{tabular}
\end{table}

These values are the fixtures expected by
\texttt{tests/quantify/test\_shahran\_worked\_example.py}; any change in
sampler, RNG seed, or parameter priors must update both the methodology
version string and this table.

\begin{thebibliography}{9}
\bibitem{ap42} U.S.\ Environmental Protection Agency. \emph{AP-42:
Compilation of Air Emissions Factors}, Section 1.3 (fuel oil combustion).
\bibitem{ipcc2006} IPCC. \emph{2006 IPCC Guidelines for National
Greenhouse Gas Inventories}, Volume 2, Chapter 1, Table 1.3.
\bibitem{wooster2005} Wooster, M.J., Roberts, G., Perry, G.L.W., Kaufman,
Y.J. (2005). Retrieval of biomass combustion rates and totals from fire
radiative power observations. \emph{Journal of Geophysical Research},
110, D24311.
\bibitem{hobbs1992} Hobbs, P.V., Radke, L.F. (1992). Airborne studies of
the smoke from the Kuwait oil fires. \emph{Science}, 256(5059), 987--991.
\bibitem{cci2026} Carbon Conflict Initiative. \emph{Iran 14-day brief}
(2026).
\bibitem{ceobs} Conflict and Environment Observatory (CEOBS), incident
records 2026.
\end{thebibliography}

\end{document}
